\section{Methodology}
\label{sec:methodology}

\subsection{Dataset}
\label{subsec:dataset}

\subsubsection{In-Distribution Data}
\label{subsubsec:id-data}

Our in-distribution (ID) data are derived from \textsc{OpenFake}~\cite{livernoche2025openfake}, a large-scale benchmark containing real photographs and synthetic images generated by multiple text-to-image models. We construct an ID subset focused on three source categories: real images, Midjourney-generated images, and Stable Diffusion-generated images. The real images in this subset are sourced from LAION, while the synthetic images are generated by Midjourney and Stable Diffusion variants.

The resulting ID dataset contains 8,989 images: 3,000 real images, 2,989 Midjourney-generated images, and 3,000 Stable Diffusion-generated images. Although the dataset is approximately balanced across these three source categories, it is imbalanced for the binary real-versus-fake classification task, since both Midjourney and Stable Diffusion images are assigned to the fake class. The binary class distribution is therefore 3,000 real images and 5,989 fake images.

The ID dataset is divided into training, validation, and test partitions using an 80/10/10 split with a fixed random seed of 42. The split is performed at the binary-label level to preserve the real-to-fake class distribution across partitions. The training set is used for model fine-tuning, the validation set is used for model selection, and the ID test set is used for final evaluation on generator families observed during training.

\begin{table}[t]
\centering
\caption{In-distribution train, validation, and test splits.}
\label{tab:id-split-counts}
\begin{tabular}{l r r r}
\hline
Split & Real & Fake & Total \\
\hline
Train & 2,400 & 4,791 & 7,191 \\
Validation & 300 & 598 & 898 \\
Test & 300 & 600 & 900 \\
\hline
Total & 3,000 & 5,989 & 8,989 \\
\hline
\end{tabular}
\end{table}

All images are resized to $256 \times 256$ pixels before model-specific preprocessing is applied.

\subsubsection{Out-of-Distribution Evaluation Data}
\label{subsubsec:ood-data}

To evaluate generalization beyond the training distribution, we construct a separate out-of-distribution (OOD) evaluation set from \textsc{OpenFake}. The OOD set is not used for training, validation, or model selection. It is used only for final evaluation.

The OOD set introduces two distribution shifts relative to the ID data. First, the fake images are generated by models that are not included in the ID train, validation, or test partitions. Second, the real images are drawn from real-image subsets distinct from those used in the ID data. This allows the OOD evaluation to measure generalization under both unseen-generator shift and real-image distribution shift.

The final OOD evaluation set contains 1,350 images: 450 real images and 900 fake images. The fake subset consists of images from 18 unseen generators, with 50 images sampled from each generator. This preserves the same approximate 1:2 real-to-fake ratio used in the ID partitions.

\begin{table}[t]
\centering
\caption{Out-of-distribution evaluation set composition.}
\label{tab:ood-composition}
\begin{tabular}{l r}
\hline
Class & Images \\
\hline
Real & 450 \\
Fake from unseen generators & 900 \\
\hline
Total & 1,350 \\
\hline
\end{tabular}
\end{table}

All OOD images are resized to $256 \times 256$ pixels before evaluation.

\subsection{Zero-Shot Evaluation}
\label{subsec:zero-shot-eval}
In the zero-shot evaluation setting, pretrained deepfake detection models were evaluated directly on the test split without additional training or domain adaptation. This experiment measures the intrinsic generalization capability of each detector against unseen diffusion-generated images.

The zero-shot setting is particularly important because modern diffusion generators follow a fundamentally different image synthesis process from GAN-based generators. Therefore, this experiment evaluates whether representations learned on older generator families transfer to diffusion-generated image artifacts.

\subsubsection{CNNDetect}
\label{subsubsec:cnndetect-zs}
For the CNN-based detector proposed by Wang et al.~\cite{wang2020cnndetect}, we used the official pretrained checkpoint \texttt{blur\_jpg\_prob0.5.pth} from the authors' public implementation.\footnote{\url{https://github.com/PeterWang512/CNNDetection}} The model was originally trained on GAN/CNN-generated images and was directly evaluated on the shared test split containing real, Midjourney-generated, and Stable Diffusion~\cite{rombach2022high}-generated images.

\subsubsection{EfficientNet}
\label{subsubsec:efficientnet-zs}
For the zero-shot evaluation of EfficientNet-B4, we used the ImageNet-pretrained checkpoint provided by the \texttt{timm} library. The model was evaluated directly on the shared test split without any additional training or domain adaptation. At inference time, the model's output logits were passed through a softmax layer to produce class probabilities, and images were assigned to the class with the higher predicted probability. No threshold tuning was applied.
\subsubsection{CLIP Zero-Shot}
\label{subsubsec:clip-zero-shot}

CLIP is a dual-encoder model: the image encoder receives images, while the text encoder receives natural-language descriptions. Although the two branches take different input types, their outputs are projected into the same image-text embedding space, which makes direct comparison possible. In the original CLIP zero-shot setting, class labels are converted into text prompts, and an image is assigned to the class whose prompt is most similar to the image embedding \cite{radford2021clip}.

We adapted this mechanism to binary real-versus-fake classification. The two candidate class descriptions were defined as \textit{``a real photograph''} and \textit{``an AI generated image''}. Since these prompts are fixed for the whole evaluation, they were encoded once using CLIP's text encoder, producing two text embeddings: one for the real class and one for the fake class. Both text embeddings lie in CLIP's shared image-text space.

For each test image, the image was preprocessed and passed through CLIP's image encoder, producing one image embedding in the same shared space. The image embedding was then compared with the two stored text embeddings. This yields two class scores per image: one for the real prompt and one for the fake prompt. The final prediction is obtained by selecting the prompt with the larger similarity score.

\subsubsection{DIRE}
\label{subsubsec:dire-zs}
For the diffusion reconstruction error detector, we used the official pretrained checkpoint lsun\_adm.pth and the ADM 256×256 unconditional diffusion model provided by the authors in the repository\footnote{\url{http://github.com/ZhendongWang6/DIRE}}. The model was evaluated directly on the test split without any additional training. Since DIRE requires reconstruction-based feature extraction prior to classification, the pipeline of creating the reconstructed image using the ADM model, and then getting the DIRE map by performing pixel wise subtraction with the original image, was applied to the test set before evaluation using the pretrained classifier

\subsubsection{FIRE}
\label{subsubsec:fire-zs}
We used the official FIRE implementation cloned from the public repository\footnote{\url{https://github.com/Chuchad/FIRE}}. We loaded the lsun\_bedroom+adm pretrained checkpoint. The model was trained on ADM-generated bedroom images from the LSUN dataset. This checkpoint is chosen for its domain mismatch with OpenFake (general-purpose photography and AI-generated art), making it the appropriate configuration for a zero-shot evaluation. Additionally, this pretrained model is used to ensure the comparison fairness between the fire and dire zeroshot results.

\subsection{Fine-Tuning}
\label{subsec:fine-tuning}

\subsubsection{Shared Training Protocol}
\label{subsubsec:shared-protocol}
To adapt pretrained detectors to diffusion-generated images, all models were fine-tuned using the shared 80/10/10 train-validation-test split. In all cases, model selection was based on validation AUC, and early stopping was applied when validation performance ceased to improve. All experiments were conducted on NVIDIA T4 GPUs provided through Google Colab.

Since the evaluated models possess significantly different architectures and parameter counts, enforcing identical layer-level fine-tuning configurations would not produce a meaningful comparison. Instead, each model was fine-tuned using an architecture-specific partial adaptation strategy, with the number and percentage of trainable parameters reported explicitly in Table~\ref{tab:ft-summary}. The optimizer, learning rate schedule, and batch size were configured independently per model to match each architecture's requirements; these are detailed in the model-specific subsections below.

Because the binary dataset contains more fake images than real images, some fine-tuning runs used a weighted loss to account for class imbalance. For each model, we report whether the loss was weighted or unweighted in Table~\ref{tab:finetuning-loss-config}.

\begin{table}[t]
\centering
\caption{Fine-tuning loss configuration for each model.}
\label{tab:finetuning-loss-config}
\begin{tabular}{l l}
\hline
Model & Loss configuration \\
\hline
CNNDetect & Unweighted cross-entropy \\
EfficientNet-B4 & Weighted cross-entropy \\
CLIP+LoRA & Unweighted cross-entropy \\
DIRE & Weighted cross-entropy \\
FIRE & Weighted composite loss \\
\hline
\end{tabular}
\end{table}

\subsubsection{CNNDetect}
\label{subsubsec:cnndetect-ft}
For the CNN-based detector proposed by Wang et al.~\cite{wang2020cnndetect}, the model was initialized using the official pretrained checkpoint \texttt{blur\_jpg\_prob0.5.pth}. During fine-tuning, the earlier convolutional layers were frozen while \texttt{layer4} and the final fully connected classification layer were unfrozen and updated.

In total, 14,966,785 out of 23,510,081 parameters were trainable, corresponding to 63.66\% of the network parameters.

Fine-tuning was performed using the Adam optimizer with a learning rate of $1\times10^{-4}$ and weight decay of $1\times10^{-4}$. The model was trained with a batch size of 8 for a maximum of 10 epochs. A \texttt{ReduceLROnPlateau} scheduler was used to reduce the learning rate when validation AUC stopped improving. Model selection was based on the highest validation AUC, and early stopping was applied with a patience of 3 epochs.

\subsubsection{EfficientNet}
\label{subsubsec:efficientnet-ft}
We selected EfficientNet-B4~\cite{tan2019efficientnet} as our backbone. Although EfficientNet-B7 achieves higher reported accuracy, it exceeded the memory capacity of our hardware (NVIDIA T4 GPUs), making EfficientNet-B4 a practical alternative.

In an initial fine-tuning trial, we adopted a two-stage training strategy in which all model parameters were frozen for the first three epochs, followed by partial unfreezing for the remaining epochs.

The trainable components were restricted to the following blocks: bn1.weight, bn1.bias, blocks.1, blocks.2, blocks.4, bn2.bias, bn2.weight, conv\_head, and classifier, amounting to 3,462,470 trainable parameters out of a total of 17,552,202 (19.7\%).

The backbone and classification head were optimized using separate learning rates to account for the difference in representational maturity between the pretrained convolutional layers and the newly initialized classifier. The backbone was updated with a learning rate of $1 \times 10^{-4}$, while the classification head was updated with a higher learning rate of $1 \times 10^{-3}$. Both parameter groups used the Adam optimizer with a weight decay of $1 \times 10^{-4}$. Training ran for a maximum of 30 epochs with a batch size of 32. Early stopping was applied with a patience of 5 epochs based on validation AUC, and the checkpoint with the highest validation AUC was retained as the final model.

\subsubsection{CLIP+LoRA}
\label{subsubsec:clip-lora-ft}

For the CLIP+LoRA model, we used \texttt{openai/clip-vit-large-patch14} as the pretrained backbone. Only the visual branch of CLIP was used for this detector. Each input image was first preprocessed according to the CLIP image preprocessing pipeline and then passed through the CLIP visual encoder. The resulting image representation was passed to a binary classification head that outputs two logits, corresponding to the real and fake classes.

LoRA adapters were applied to the attention layers of the CLIP visual transformer. In particular, adapters were inserted into the query and value projection matrices, while the original pretrained CLIP weights remained frozen. During training, only the LoRA adapter weights and the classification head were updated. This allows the model to adapt the CLIP visual representation to the real-versus-fake task while modifying only a small fraction of the full model.

The classification head maps the CLIP image representation to the binary output space. The complete fine-tuning configuration is summarized in Table~\ref{tab:clip_lora_config}.

\begin{table}[h]
\centering
\caption{CLIP+LoRA fine-tuning configuration.}
\label{tab:clip_lora_config}
\begin{tabular}{ll}
\hline
\textbf{Category} & \textbf{Configuration} \\
\hline
Backbone & \texttt{openai/clip-vit-large-patch14} \\
Input branch & Visual branch only \\
Trainable modules & LoRA adapters, classification head \\
LoRA target modules & Query and value projections \\
LoRA hyperparameters & $r=16$, $\alpha=32$, dropout $=0.1$ \\
Classification head & 768 $\rightarrow$ 256 $\rightarrow$ 2 \\
Trainable parameters & 1.77M / 305.74M (0.58\%) \\
Loss & Unweighted Cross-entropy \\
Optimizer & AdamW \\
Learning rate & $2 \times 10^{-4}$ \\
Weight decay & 0.01 \\
Batch size & 32 \\
Scheduler & ReduceLROnPlateau on validation AUC \\
Model selection & Best validation AUC \\
Early stopping & Validation AUC, patience = 1 epoch \\
\hline
\end{tabular}
\end{table}


\subsubsection{DIRE}
\label{subsubsec:dire-ft}
The first step was applying the DIRE pipeline to the full dataset to compute DIRE maps for all training and validation splits before fine-tuning began. Crucially, the test DIRE maps computed during the zero-shot phase were reused without recomputation to ensure consistency across both evaluation phases.
Due to the significant computational cost of running the ADM model for each image, the DIRE map generation was implemented with a resumable pipeline that persisted computed maps to Google Drive, allowing computation to be distributed across multiple Colab sessions (approximately 20 hours total).
The DIRE map generation pipeline (ADM inversion and reconstruction) remained unchanged, only the ResNet-50 classifier was fine-tuned. During training, only the final bottleneck block of \texttt{layer4} and the fully connected classification head were unfrozen, constituting approximately 4,466,690 trainable parameters out of 23,512,130 total ($\approx$19\%).

A WeightedRandomSampler was applied during training to address the class imbalance present in the dataset.

Training followed the shared experimental protocol: Adam optimizer with learning rate 1e-4, weight decay 1e-4, CosineAnnealingLR with ($T_{\max}=10$), batch size 8, and 10 epochs with the model being saved at the best checkpoint selected by validation AUC.

\subsubsection{FIRE}
\label{subsubsec:fire-ft}
For the FIRE model, the same pretrained ADM backbone (\texttt{Lsunbedroom\_adm}) used for DIRE was adopted to maintain consistency and ensure that performance differences can be attributed to the detection methodology rather than the diffusion model.
The FIRE architecture consists of three main submodules: a Variational Autoencoder (VAE) containing 83,653,863 parameters, a ResNet-based feature extraction module with 23,466,369 parameters, and a Fast Fourier Transform (FFT) filter module containing 23,912 parameters, totaling 107,144,144 parameters.
During fine-tuning, only 22,046,057 parameters (20.6\%) were updated: specifically, the higher-level feature extraction layers (layer~3 and layer~4), the fully connected classification layer, and the FFT filter module. The VAE and lower ResNet layers remained frozen.

The FIRE training objective combines binary classification and reconstruction-based supervision using a weighted composite loss:
\begin{equation}
\mathcal{L}_{\text{FIRE}} = 0.6 \times \mathcal{L}_{\text{cross-entropy}} + 0.2 \times \mathcal{L}_{\text{MSE-recon}} + 0.2 \times \mathcal{L}_{\text{MSE-mask}}
\end{equation}
where $\mathcal{L}_{\text{cross-entropy}}$ is the binary classification loss, $\mathcal{L}_{\text{MSE-recon}}$ is the reconstruction consistency loss, and $\mathcal{L}_{\text{MSE-mask}}$ penalizes reconstruction differences within masked (mid-frequency filtered) regions. This composite loss follows the original FIRE implementation.

The model was trained using Adam with learning rate $1 \times 10^{-4}$, weight decay $1 \times 10^{-4}$, a cosine annealing learning rate scheduler ($T_{\max}=10$), and batch size~8. Training ran for a maximum of 25 epochs with early stopping (patience~5) based on validation AUC. The best model was obtained at epoch~16 (approximately 12 hours of training).

\subsection{Summary of Fine-Tuning Configurations}
\label{subsec:ft-summary}

Table~\ref{tab:ft-summary} summarizes the fine-tuning configuration for each model to facilitate comparison.

\begin{table}[H]
\centering
\caption{Fine-Tuning Configuration Summary}
\label{tab:ft-summary}
\begin{tabular}{|l|c|c|c|}
\hline
\textbf{Model} & \textbf{Total Params} & \textbf{Trainable} & \textbf{\% Trainable} \\
\hline
CNNDetect & 23.5M & 15.0M & 63.7\% \\
\hline
EfficientNet-B4 & 17.6M & 3.5M & 19.7\% \\
\hline
CLIP+LoRA & 305.74M & 1.77M & 0.58\% \\
\hline
DIRE (classifier) & 23.5M & 4.5M & 19.0\% \\
\hline
FIRE & 107.1M & 22.0M & 20.6\% \\
\hline
\end{tabular}
\end{table}
